\prelimchapter{Introducción}{\topbanner}\label{cap:introduccion}
\vspace*{0.7cm}


En la actualidad, las organizaciones dependen cada vez más de servicios y recursos digitales para soportar sus actividades académicas, administrativas y de investigación. Este crecimiento tecnológico ha impulsado la necesidad de implementar mecanismos que permitan administrar de manera eficiente las identidades digitales, los accesos y las credenciales de los usuarios dentro de las infraestructuras informáticas. Según \citet{Pohn2022}, la gestión de identidades y accesos es un componente crucial de la infraestructura de \ac{TI}, ya que mantiene atributos, credenciales, roles y permisos asociados a un identificador, y busca ofrecer una visión común e interoperabilidad entre distintos enfoques de identidad.

Los servicios de directorio desempeñan un papel importante dentro de este contexto, puesto que permiten organizar y administrar información asociada a usuarios, grupos, recursos y políticas de acceso dentro de una infraestructura tecnológica. En términos generales, un servicio de directorio corresponde a una estructura jerárquica destinada al almacenamiento y consulta centralizada de información relacionada con objetos dentro de una red \citep{Microsoft2025}. Asimismo, el protocolo \ac{LDAP} proporciona un mecanismo estandarizado para acceder y mantener servicios de directorio, permitiendo que múltiples sistemas y aplicaciones consulten de manera centralizada la información de identidad y autenticación de los usuarios \citep{Pohn2022}.

En instituciones académicas y centros de investigación, la administración independiente de usuarios en múltiples plataformas puede generar dificultades relacionadas con el control de accesos, la duplicidad de credenciales y la aplicación uniforme de políticas de seguridad. En este contexto, la gestión de identidades, el control de acceso y la autorización constituyen procesos fundamentales para administrar perfiles de usuario, asignar roles y privilegios, y garantizar que únicamente los usuarios autorizados accedan a los recursos institucionales. Según \citet{Mollakuqe2024}, estas prácticas desempeñan un papel esencial en la protección de los sistemas y recursos digitales dentro de las instituciones de educación superior.

El \ac{GRID} de la Universidad del Quindío dispone de diversos servicios tecnológicos que soportan actividades académicas, investigativas y administrativas. Sin embargo, el crecimiento progresivo de esta infraestructura ha implicado la incorporación de múltiples plataformas que actualmente administran usuarios y credenciales de forma independiente. Esta situación dificulta la gestión centralizada de accesos, incrementa los procesos manuales de administración y limita la implementación uniforme de controles de seguridad y auditoría.

De acuerdo con \citet{Felts2023}, la autenticación federada simplifica el acceso a recursos digitales al permitir que los usuarios utilicen las credenciales de su institución de origen, además de mejorar la experiencia de acceso y reducir esfuerzos operativos asociados a la administración de los mecanismos de autenticación. Asimismo, \citet{Prout2019} señalan que la autenticación federada puede reducir significativamente la carga de mantenimiento de cuentas, mejorar la experiencia de los usuarios y fortalecer la seguridad al aprovechar los sistemas de gestión de identidad y los mecanismos de autenticación implementados por las organizaciones de origen.

Por otra parte, la norma ISO/IEC 27001:2022 establece requisitos para implementar, mantener y mejorar continuamente un sistema de gestión de seguridad de la información, incluyendo controles relacionados con la gestión de accesos y la protección de los activos de información \citep{ISO27001}. Asimismo, la norma ISO/IEC 25010:2011 define un modelo de calidad para sistemas y productos software que contempla características como seguridad, compatibilidad, fiabilidad, mantenibilidad y adecuación funcional, aspectos relevantes en la implementación de soluciones orientadas a la gestión de identidades y accesos \citep{ISO25010}.

En este contexto, el presente proyecto tiene como propósito implementar un servicio de directorio para la infraestructura tecnológica del grupo de investigación \ac{GRID}, orientado a centralizar la gestión de usuarios, autenticación y control de accesos mediante el uso de tecnologías de código abierto. Para ello, se plantea un proceso compuesto por la identificación de necesidades funcionales y técnicas, la caracterización y evaluación de alternativas tecnológicas mediante el método \ac{DAR}, el diseño arquitectónico de la solución seleccionada y la implementación de un prototipo funcional adaptado al entorno del \ac{GRID}.

Finalmente, el desarrollo de este proyecto busca contribuir al fortalecimiento de la seguridad, la organización y la administración eficiente de los recursos tecnológicos del grupo de investigación \ac{GRID}, estableciendo una base centralizada y escalable para futuras integraciones y mecanismos avanzados de autenticación y gestión de identidades.
